\documentclass[10pt,a4paper]{amsart}
\usepackage{amsmath,amssymb,amsthm,graphicx}
\usepackage{verbatim}
\newenvironment{reference}{\comment}{\endcomment}
%\newenvironment{reference}{}{}
\newenvironment{slogan}{\comment}{\endcomment}
\newenvironment{history}{\comment}{\endcomment}

% For commutative diagrams we use Xy-pic
\usepackage[all]{xy}

% We use 2cell for 2-commutative diagrams.
\xyoption{2cell}
% Theorem environments.
%
\theoremstyle{plain}
\newtheorem{theorem}[subsection]{Theorem}
\newtheorem{proposition}[subsection]{Proposition}
\newtheorem{lemma}[subsection]{Lemma}

\theoremstyle{definition}
\newtheorem{definition}[subsection]{Definition}
\newtheorem{example}[subsection]{Example}
\newtheorem{exercise}[subsection]{Exercise}
\newtheorem{situation}[subsection]{Situation}

\theoremstyle{remark}
\newtheorem{remark}[subsection]{Remark}
\newtheorem{remarks}[subsection]{Remarks}

\numberwithin{equation}{subsection}

% Macros
%
\def\lim{\mathop{\mathrm{lim}}\nolimits}
\def\colim{\mathop{\mathrm{colim}}\nolimits}
\def\Spec{\mathop{\mathrm{Spec}}}
\def\Hom{\mathop{\mathrm{Hom}}\nolimits}
\def\Ext{\mathop{\mathrm{Ext}}\nolimits}
\def\SheafHom{\mathop{\mathcal{H}\!\mathit{om}}\nolimits}
\def\SheafExt{\mathop{\mathcal{E}\!\mathit{xt}}\nolimits}
\def\Sch{\mathit{Sch}}
\def\Mor{\mathop{\mathrm{Mor}}\nolimits}
\def\Ob{\mathop{\mathrm{Ob}}\nolimits}
\def\Sh{\mathop{\mathit{Sh}}\nolimits}
\def\NL{\mathop{N\!L}\nolimits}
\def\CH{\mathop{\mathrm{CH}}\nolimits}
\def\proetale{{pro\text{-}\acute{e}tale}}
\def\etale{{\acute{e}tale}}
\def\QCoh{\mathit{QCoh}}
\def\Ker{\mathop{\mathrm{Ker}}}
\def\Im{\mathop{\mathrm{Im}}}
\def\Coker{\mathop{\mathrm{Coker}}}
\def\Coim{\mathop{\mathrm{Coim}}}

% Boxtimes
%
\DeclareMathSymbol{\boxtimes}{\mathbin}{AMSa}{"02}

%
% Macros for moduli stacks/spaces
%
\def\QCohstack{\mathcal{QC}\!\mathit{oh}}
\def\Cohstack{\mathcal{C}\!\mathit{oh}}
\def\Spacesstack{\mathcal{S}\!\mathit{paces}}
\def\Quotfunctor{\mathrm{Quot}}
\def\Hilbfunctor{\mathrm{Hilb}}
\def\Curvesstack{\mathcal{C}\!\mathit{urves}}
\def\Polarizedstack{\mathcal{P}\!\mathit{olarized}}
\def\Complexesstack{\mathcal{C}\!\mathit{omplexes}}
% \Pic is the operator that assigns to X its picard group, usage \Pic(X)
% \Picardstack_{X/B} denotes the Picard stack of X over B
% \Picardfunctor_{X/B} denotes the Picard functor of X over B
\def\Pic{\mathop{\mathrm{Pic}}\nolimits}
\def\Picardstack{\mathcal{P}\!\mathit{ic}}
\def\Picardfunctor{\mathrm{Pic}}
\def\Deformationcategory{\mathcal{D}\!\mathit{ef}}
\usepackage{cleveref}
\crefname{originalequation}{original equation}{original equations}
\Crefname{originalequation}{Original equation}{Original equations}
\makeatletter
\newlabel{equation-multiplicity}{{0BI6}{source}{}{sourcepointer-1}{}}
\newlabel{equation-multiplicity@cref}{{[originalequation][0][]0BI6}{[1][0][]source}{}{}{}}
\makeatother
\crefname{originalreference}{original reference}{original references}
\Crefname{originalreference}{Original reference}{Original references}
\makeatletter
\newlabel{lemma-blowup-curve}{{0BI5}{source}{}{sourcepointer-2}{}}
\newlabel{lemma-blowup-curve@cref}{{[originalreference][0][]0BI5}{[1][0][]source}{}{}{}}
\makeatother
\crefname{originalreference}{original reference}{original references}
\Crefname{originalreference}{Original reference}{Original references}
\makeatletter
\newlabel{lemma-blowup-nonsingular-curves-meeting-at-point}{{0BI7}{source}{}{sourcepointer-3}{}}
\newlabel{lemma-blowup-nonsingular-curves-meeting-at-point@cref}{{[originalreference][0][]0BI7}{[1][0][]source}{}{}{}}
\makeatother
\crefname{originalreference}{original reference}{original references}
\Crefname{originalreference}{Original reference}{Original references}
\makeatletter
\newlabel{lemma-resolve-curve}{{0BI4}{source}{}{sourcepointer-4}{}}
\newlabel{lemma-resolve-curve@cref}{{[originalreference][0][]0BI4}{[1][0][]source}{}{}{}}
\makeatother
\title{Separating and smoothing finitely many curves on a Noetherian surface by blowups in points}
\author{The Stacks Project Authors}
\date{stacks-project:resolve.tex:local-snapshot-20261012; GNU FDL 1.2 or later}
\begin{document}
\maketitle
\noindent\textit{Editorial note.} All delivered mathematical text is the author's own, copied from the pinned Stacks Project chapter file resolve.tex; no statement or proof step was rewritten, reordered or removed.\par\medskip
\noindent\textit{Editorial note.} The counted proof is the author's complete argument: the reduction to regular curves, the decreasing maximal intersection multiplicity, and the separation step when the maximum has reached one.\par\medskip
\noindent\textit{Editorial note.} The label lemma-blowup-curves names the counted statement, and the proof environment immediately following it in the source is delivered verbatim as its proof.\par\medskip
\noindent\textit{Editorial note.} The counted statement is carried by the author's own lemma environment in the chapter source.\par\medskip
\noindent\textit{Editorial note.} Pre-open gate: python operations/agent-authoring/tools/precard\_gate.py --source resolve.tex --label lemma-blowup-curves --start 4394 --end 4417 --local-id 28119 returned RESULT=OK (L28119 free, occupancy FREE, public interval clear). Pre-assembly tag self-check: canonical tag resolved and not present in reference/public-items.jsonl.\par\medskip
\noindent\textit{Editorial note.} The statement and proof are copied from the pinned chapter file; the cited inputs are kept as external pointers with their canonical Stacks tags.\par\medskip
\noindent\textit{Allowed standard prerequisites (editorial).}\begin{itemize}
\item Making each curve regular by blowups in closed points. Same chapter, the lemma resolving an integral curve on a Noetherian surface by blowups in closed points. Reduces to the case where all the curves are already regular, which is the hypothesis of the separating argument.
\item Descent of the intersection multiplicity when blowing up a common point of two regular curves. Same chapter, the computation of the behaviour of the intersection multiplicity of two curves meeting at a point under blowing up that point. Provides the decreasing invariant used in the induction and the separation of curves meeting with multiplicity one.
\end{itemize}\medskip
\ifdefined\hypertarget\hypertarget{sourcepointer-1}{}\fi
\noindent\textit{Source reference.} Original equation 0BI6: sources/stacks-resolve/pinned-original-source.tex, source lines 4333--4333. Prerequisite pointer invoked at proof line(s) 4411; it names a standard fact of the same chapter that the counted argument uses, and every step of the counted proof built on it is retained verbatim in the delivered body.\par\medskip
\ifdefined\hypertarget\hypertarget{sourcepointer-2}{}\fi
\noindent\textit{Source reference.} Original reference 0BI5: sources/stacks-resolve/pinned-original-source.tex, source lines 4307--4307. Prerequisite pointer invoked at proof line(s) 4408; it names a standard fact of the same chapter that the counted argument uses, and every step of the counted proof built on it is retained verbatim in the delivered body.\par\medskip
\ifdefined\hypertarget\hypertarget{sourcepointer-3}{}\fi
\noindent\textit{Source reference.} Original reference 0BI7: sources/stacks-resolve/pinned-original-source.tex, source lines 4346--4346. Prerequisite pointer invoked at proof line(s) 4413; it names a standard fact of the same chapter that the counted argument uses, and every step of the counted proof built on it is retained verbatim in the delivered body.\par\medskip
\ifdefined\hypertarget\hypertarget{sourcepointer-4}{}\fi
\noindent\textit{Source reference.} Original reference 0BI4: sources/stacks-resolve/pinned-original-source.tex, source lines 4248--4248. Prerequisite pointer invoked at proof line(s) 4398; it names a standard fact of the same chapter that the counted argument uses, and every step of the counted proof built on it is retained verbatim in the delivered body.\par\medskip
\begin{lemma}
\label{lemma-blowup-curves}
Let $X$ be a Noetherian scheme. Let $Y_i \subset X$, $i = 1, \ldots, n$
be an integral closed subschemes of dimension $1$ each satisfying the
equivalent conditions of Lemma \ref{lemma-resolve-curve}. Then there
exists a finite sequence
$$
X_n \to X_{n - 1} \to \ldots \to X_1 \to X
$$
of blowups in closed points such that the strict transform $Y'_i \subset X_n$
of $Y_i$ in $X_n$ are pairwise disjoint regular curves.
\end{lemma}

\begin{proof}
It follows from Lemma \ref{lemma-blowup-curve} that we may assume $Y_i$
is a regular curve for $i = 1, \ldots, n$. For every $i \not = j$
and $p \in Y_i \cap Y_j$ we have the invariant
$m_p(Y_i \cap Y_j)$ (\ref{equation-multiplicity}).
If the maximum of these numbers is $> 1$, then we can decrease
it (Lemma \ref{lemma-blowup-nonsingular-curves-meeting-at-point})
by blowing up in all the points $p$ where the maximum is attained.
If the maximum is $1$ then we can separate the curves using the
same lemma by blowing up in all these points $p$.
\end{proof}

\end{document}
